% common-scrbook/preamble/setup.tex
%
% Base preamble for the scrbook target. Loaded by main.tex after
% \documentclass{scrbook}. Establishes language, math, fonts, hyperlinks,
% cross-references, and color palette. Segment-specific commands
% (\segmenthead, \eqtag, \segmentfield) and environments (workingnotes,
% callout) live in segment.tex / environments.tex.

%----------------------------------------------------------------------------
% Language & quoting
%----------------------------------------------------------------------------
\usepackage[english]{babel}
\usepackage[english=american]{csquotes}

%----------------------------------------------------------------------------
% Math — amsmath family, mathtools (which patches several amsmath rough
% edges), then the unicode-math layer for fontspec-driven math fonts.
%----------------------------------------------------------------------------
\usepackage{amsmath}
\usepackage{amssymb}
\usepackage{mathtools}
\usepackage{amsthm}
\usepackage{unicode-math}

%----------------------------------------------------------------------------
% Fonts.
%
% Default is STIX Two — a single coordinated text + math family that ships
% with current TeX Live and renders with the visual register of a classical
% mathematical monograph (the AMS GSM series, Folland's Real Analysis,
% Lang's Algebra). Both faces are open-source and machine-portable.
%
% To swap fonts: replace the three blocks below. fontspec finds OS fonts via
% fontconfig and TeX Live fonts via its texmf tree; either source works.
%----------------------------------------------------------------------------
\setmainfont{STIX Two Text}[
  Ligatures = {TeX, Common},
]
\setmathfont{STIX Two Math}

% Sans-serif: only used by KOMA scrbook for chapter / section / subsection
% headings. Latin Modern Sans is the safe portable default; matches the
% serif body weight reasonably well without pulling in a third family.
\setsansfont{Latin Modern Sans}[
  Ligatures = TeX,
  Scale     = MatchLowercase,
]

% Monospace: Latin Modern Mono — narrow enough to sit comfortably in inline
% code spans (slug names, file paths) without crowding surrounding prose.
\setmonofont{Latin Modern Mono}[
  Scale = MatchLowercase,
]

%----------------------------------------------------------------------------
% GFM/Obsidian math-compat macros — same shims as the kaobook target.
% Source segments use \lt and \gt because raw < and > break GitHub's math
% renderer (interpreted as HTML tags). LaTeX math knows neither — define
% them as the obvious operators.
%----------------------------------------------------------------------------
\providecommand{\lt}{<}
\providecommand{\gt}{>}

%----------------------------------------------------------------------------
% Microtypography
%----------------------------------------------------------------------------
\usepackage{microtype}

%----------------------------------------------------------------------------
% Widow / orphan suppression. Same discipline as the kaobook target — a
% stranded single line at top or bottom of a page reads as broken; the
% occasional extra whitespace is the right trade.
%----------------------------------------------------------------------------
\widowpenalty=10000
\clubpenalty=10000

%----------------------------------------------------------------------------
% String predicates — used by segment.tex for the eq-tag dispatch and any
% conditional rendering.
%----------------------------------------------------------------------------
\usepackage{xstring}

%----------------------------------------------------------------------------
% Color palette. Modest by design — link/cross-ref navy, a quiet rule-gray,
% a muted sand for the working-notes register. No status-tier colors:
% scrbook target drops the status badges entirely.
%----------------------------------------------------------------------------
\usepackage{xcolor}
\definecolor{rule-gray}      {HTML}{B8B8B8}
\definecolor{link-navy}      {HTML}{1F3A6F}
\definecolor{wn-rule}        {HTML}{8C8676}
\definecolor{wn-fg}          {HTML}{6E6354}
\definecolor{callout-rule}   {HTML}{6B7B95}

%----------------------------------------------------------------------------
% pdfpages — for the cover page (rsvg-convert produces a single-page PDF
% that we \includepdf at the very front).
%----------------------------------------------------------------------------
\usepackage{pdfpages}

%----------------------------------------------------------------------------
% needspace — \needspace{N\baselineskip} prevents a heading from being
% orphaned at page bottom by forcing a page break first when N lines aren't
% available below.
%----------------------------------------------------------------------------
\usepackage{needspace}

%----------------------------------------------------------------------------
% Tables — same toolkit as the kaobook target so AsfLatex's table converter
% (in segment_renderer.rb) renders identically here.
%----------------------------------------------------------------------------
\usepackage{tabularx}
\usepackage{booktabs}
\usepackage{array}
\usepackage{caption}

%----------------------------------------------------------------------------
% Equation + table numbering reset per chapter — same convention as the
% kaobook target (3.5, 3.6, …, then 4.1 in next chapter).
%----------------------------------------------------------------------------
\numberwithin{equation}{chapter}
\numberwithin{table}{chapter}

%----------------------------------------------------------------------------
% Hyperlinks + cross-references. hyperref must load late (it patches many
% other packages). cleveref must load AFTER hyperref so its \cref renders
% as a clickable link.
%----------------------------------------------------------------------------
\usepackage[
  colorlinks = true,
  linkcolor  = link-navy,
  citecolor  = link-navy,
  urlcolor   = link-navy,
  filecolor  = link-navy,
  pdfborder  = {0 0 0},
]{hyperref}

\usepackage[capitalize,noabbrev]{cleveref}

% Cross-ref display: bare number for sections and chapters. Resolved cross-
% refs in the assembled markdown look like `[Definition 1.4](#def-foo)`;
% the converter emits `\cref{seg:def-foo}`. We want LaTeX to produce the
% number "1.4" — the type word ("Definition") was kept in the markdown's
% visible link text and survives in source prose like "see Definition
% \cref{seg:def-foo}".
%
% IMPORTANT: cleveref's format placeholder for the stored label value is
% `#1`, NOT `\thesection`. `#2#1#3` is the canonical
% "hyperref-open + value + hyperref-close" idiom.
\AtBeginDocument{%
  \crefformat{section}{#2#1#3}%
  \Crefformat{section}{#2#1#3}%
  \crefrangeformat{section}{#3#1#4--#5#2#6}%
  \Crefrangeformat{section}{#3#1#4--#5#2#6}%
  \crefmultiformat{section}{#2#1#3}{, #2#1#3}{, #2#1#3}{, #2#1#3}%
  \Crefmultiformat{section}{#2#1#3}{, #2#1#3}{, #2#1#3}{, #2#1#3}%
  % Appendix segments are chapters; mirror the section convention so
  % chapter and section cross-refs read identically in running text.
  \crefformat{chapter}{#2#1#3}%
  \Crefformat{chapter}{#2#1#3}%
  \crefrangeformat{chapter}{#3#1#4--#5#2#6}%
  \Crefrangeformat{chapter}{#3#1#4--#5#2#6}%
  \crefmultiformat{chapter}{#2#1#3}{, #2#1#3}{, #2#1#3}{, #2#1#3}%
  \Crefmultiformat{chapter}{#2#1#3}{, #2#1#3}{, #2#1#3}{, #2#1#3}%
}

%----------------------------------------------------------------------------
% \ifdefempty — used by main.tex to detect empty cover-path. Provided by
% etoolbox, which KOMA pulls in transitively, but load explicitly so the
% conditional always resolves.
%----------------------------------------------------------------------------
\usepackage{etoolbox}
